\chapter{Conclusion and Future Work}
\label{ch:conclusion}

\section{What was established}

This thesis set out to determine, by measurement, where a production
visual-SLAM workload's computation belongs on the emerging spectrum of
memory-centric hardware. Its answers, in the order of the research
questions:

\textbf{RQ1 --- the methodology exists.} GPU-DAMOV re-derives every
component of the reference data-movement methodology for GPUs---screening,
machine-independent locality, threshold classification, and the full
robustness battery---with the three GPU-forced redesigns (two-level LFMR,
occupancy-conditioned latency classes, a coalescing axis) and with DAMOV's
simulated intervention replaced by measured ones. The parity argument is
componentwise and complete at DAMOV's granularity
(\S\ref{sec:met-parity}).

\textbf{RQ2 --- the classes are real and stable.} Eight bottleneck classes
(G0--G7, one of them---dependency-bound G7---forced into existence by the
data) classify 49 kernels with 91\% cross-sequence consistency, survive
$\pm$25\% threshold stress, are recovered by two label-free clustering
algorithms (purity $\approx$0.68 both), agree 80\% across two GPU
microarchitectures, and correctly classify a known-truth calibration suite
after three archetype defects were fixed with the classifier untouched.

\textbf{RQ3 --- every byte has a name.} The TaggedAllocator attribution
resolves all 274 allocations and gives 48/49 kernels a unanimous
data-structure tag across the corpus. Its two decisive findings: the GPU
memory budget is \emph{static} (the growing keyframe database lives
host-side, its GPU-resident state a fixed 6.7\,MB), and 92.6\% of the
loop-closure scan's DRAM volume is register-spill scratch---a compiler
artifact no counter-only study could have separated from data traffic.

\textbf{RQ4 --- the placement.} The workload splits three ways: a
cache-immune streaming front-end (near-sensor placement; no cache size
helps, measured); a hot-persistent solver that should stay on the GPU at
this scale; and a cold-persistent, session-growing database tier whose scan
is a genuine scattered gather (in/near-storage plus scatter-capable PiM
candidacy). Time-weighted, 61--72\% of GPU time carries near-data affinity;
16--26 of 49 kernels clear an analytical offload bar under conservative-to-
moderate scenarios; and the baselines a follow-on study must beat are now
measured: 34.67\,J whole-run energy, 1.68\,MB/frame sensor upload, 41\% of
kernel time in boundary copies, 66\,MB storage ingestion, 708\,MB peak host
memory.

\textbf{RQ5 --- the measurements are trustworthy.} The characterized system
reproduces its vendor's published accuracy (141-configuration matrix);
profiling is output-neutral (bit-identical trajectories on deterministic
modes across 192 variants); the noise floor is measured (CoV 0.14\%); and
the project's two self-corrections---most notably the reversal of an
interim ``coalesced streaming'' conclusion once register-spill traffic was
separated from data---are documented in place, ending with independent
instruments in agreement.

\section{The thesis, restated}

For Physical-AI perception workloads, the question ``should we build PiM?''
is ill-posed; the well-posed question is ``\emph{which data structure},
touched by \emph{which kernel}, under \emph{which access pattern}, belongs
on \emph{which substrate}?''---and it is answerable today, on production
software, with commodity instruments, to the granularity of a named
allocation and with quantified confidence. Answered for cuVSLAM, it yields
not one accelerator but a placement: sensor-side consumption for streams,
the GPU for the solver, near-storage and scatter-capable near-memory
engines for the map---with the measured evidence attached to every cell.

\section{Future work}
\label{sec:roadmap}

The project's phased roadmap continues past this document:

\begin{enumerate}
\item \textbf{Complete the validation suite.} Execute the scripted
clock-domain intervention (per-class predictions are stated ex ante:
bandwidth classes track the memory clock, compute/on-chip/dependency classes
the core clock) and fold its verdicts into the artifact.
\item \textbf{Name the last bytes.} The kernel-argument correlation layer
(Layer 3) resolves the two named residual kernels carrying 83\% of unmapped
traffic; its target list and join interface are already committed.
\item \textbf{Embedded-target re-derivation.} Re-run the regime on
Jetson-class hardware (the harness already drives remote targets) to
re-derive the codegen-dependent quantities---above all the spill share---and
test the taxonomy's transfer to the deployment-class part.
\item \textbf{The architecture study.} Feed the traces to Accel-Sim under
the generated NDP configurations~\cite{khairy2020accelsim}, with
AccelWattch~\cite{kandiah2021accelwattch} for per-class energy, reporting
speedup and energy \emph{deltas} against this thesis's measured baselines
---per the standing rule, never absolutes. The conservative/moderate
placement scenarios of \S\ref{sec:syn-model} become simulated designs; the
16--26-kernel offload sets are the experiment list.
\item \textbf{Breadth via adapters.} Run non-SLAM GPU workloads (DNN
inference, GPU database operators) through the unchanged pipeline to test
the methodology's generality claim---the adapter interface exists precisely
so that breadth costs execution time, not engineering.
\end{enumerate}

The characterization is complete; the design space it opens is measured,
bounded, and---for the first time for a production Physical-AI perception
stack---named, byte by byte.
